\documentclass[11pt]{article}

\usepackage[margin=1in]{geometry}
\usepackage{booktabs}
\usepackage{hyperref}
\usepackage{amsmath,amssymb}
\usepackage{enumitem}
\usepackage{microtype}
\setlength{\emergencystretch}{2em}

\title{A Data-Centric Comparison of Autoregressive, Energy-Based, and Novel Energy Models for Verifier-Guided Generation}
\author{Brando Miranda}
\date{\today \\ \emph{Planning draft}}

\begin{document}
\maketitle

\begin{abstract}
This paper compares three primary model families under one data-centric protocol:
an autoregressive language-model baseline, a conventional energy-based model,
and a novel free-energy/post-softmax energy model.
It also includes diffusion/iterative denoising as a required reference point,
because otherwise the comparison artificially excludes the strongest
non-autoregressive generative baseline.
The contribution is the controlled comparison itself: common splits, common
candidate pools, verifier-backed metrics, and matched compute reporting on toy
controls, VeriBench/Lean, and an initial MNIST vision track.
\end{abstract}

\section{Thesis}

The paper should be publishable even if the novel EBM is only competitive or
partially successful.
The scientific claim is that architecture debates are underdetermined unless the
data, verifier, candidate pool, compute budget, and evaluation protocol are held
fixed.

\section{Model families}

\begin{enumerate}[leftmargin=2em]
  \item \textbf{Autoregressive/LLM baseline.}
        A standard next-token model with teacher forcing, softmax attention, and
        softmax output head.
  \item \textbf{Normal EBM baseline.}
        A conventional energy model trained with established EBM objectives such
        as contrastive divergence, noise-contrastive estimation, or score-style
        surrogates where appropriate.
  \item \textbf{Novel EBM.}
        The proposed free-energy/post-softmax model from this project.
  \item \textbf{Diffusion / iterative denoising reference.}
        A DDPM-style or masked-diffusion-style baseline that pays iterative
        inference cost rather than local autoregressive normalization.
\end{enumerate}

\section{Data}

\begin{itemize}[leftmargin=2em]
  \item \textbf{Toy controls:} mechanisms with known ground truth.
  \item \textbf{VeriBench/Lean:} primary real-data domain with a hard verifier.
  \item \textbf{MNIST:} first vision smoke test for ordering/global-validity effects.
  \item \textbf{Later vision dataset:} to decide after the MNIST methodology is
        stable.
\end{itemize}

\section{Metrics}

\begin{center}
\begin{tabular}{ll}
\toprule
Metric & Purpose \\
\midrule
Pass@k & verifier-backed success \\
cost per verified solution & efficiency under search/retry \\
success vs length/depth & test error-compounding hypotheses \\
candidate energy/rank quality & evaluate EBM scoring \\
matched-compute loss/pass rate & prevent unfair architecture comparisons \\
\bottomrule
\end{tabular}
\end{center}

\section{Experimental protocol}

Every run reports the split, model class, parameter count, training compute,
inference budget, verifier budget, seeds, and bootstrap intervals where
applicable.
The data protocol is the paper's main contribution, so leakage boundaries and
candidate-pool construction should be described before model results.

\section{Expected outcomes}

A positive result shows that the novel EBM improves verifier-backed success or
cost per verified solution.
A useful negative result shows that the data-centric protocol exposes a precise
failure mode of EBM training, EBM inference, or AR baselines.

\end{document}
